\chapter*{Введение} % * не проставляет номер

\addcontentsline{toc}{chapter}{Введение} % вносим в содержание

Светофорные объекты (СО) - это ключевой элемент инфраструктуры на дорожных участках.

Естественно заметить, что светофоры организовывают движение транспортных средств, предоставляя им возможность перемещаться по дорогам в соответствии с определенными правилами и приоритетами. Правильно настроенные программы светофорных объектов могут сгладить поток транспортных средств, предотвратить образование заторов и улучшить общую пропускную способность участка дороги.

Соответственно, необходимо подобрать оптимальный режим работы светофора так, чтобы удовлетворить желаемым критериям, например, минимизировать время ожидания в пробке.

Нетрудно догадаться, что настраивать длительности фаз и времена действия регулирующих сигналов в реальном времени на дороге - задача труднореализуемая и опасная. Небольшие изменения в работе программ регулирования трафика могут не дать значимо отличающихся от имеющихся результатов, а радикальные изменения режимов работы создают опасность возниковения аварийных ситуаций на дороге.


Чтобы избежать вышеизложенных проблем, применяется подход, называющийся микромоделирование автомобильного трафика. На основе реальных статистических данных возможно построить правдоподобную модель ситуации на дороге. При помощи настройки параметров модели (длительностей регулирующих фаз светофоров, времён действия сигналов) можно проводить виртуальные эксперименты, анализировать воздействие различных режимов работы светофоров на поток транспорта.


В данной работе реализована микромодель дорожного трафика с использованием инструмента SUMO \cite{sumo-docs}. Настройку параметров модели предлагается производить на основании эволюционных алгоритмов: генетический алгоритм \cite{gad2023pygad}, метод роя частиц \cite{pyswarms2018doc} и эволюционная стратегия с адаптацией матрицы ковариаций \cite{HansenCmaes}. По результатам оптимизации параметров микромодели предложено улучшение режимов работы светофорных объектов.



%% Вспомогательные команды - Additional commands
%\newpage % принудительное начало с новой страницы, использовать только в конце раздела
%\clearpage % осуществляется пакетом <<placeins>> в пределах секций
%\newpage\leavevmode\thispagestyle{empty}\newpage % 100 % начало новой строки